% Shared exercise chunk: l1-rope
The following tool may be used in \cref{ex:rope}.
\begin{theorem}[Hoeffding's inequality]
\label{thm:hoeffding}
Let $n\ge1$ be an integer and let $Z_1,\ldots,Z_n$ be independent real
random variables. Suppose $Z_r\in[a_r,b_r]$ almost surely for each $r\in[n]$,
where $a_r<b_r$ are real numbers. For every real $t\ge0$,
\[
\mathbb P\!\left(\sum_{r=1}^n (Z_r-\mathbb E[Z_r])\ge t\right)
\le \exp\!\left(-\frac{2t^2}{\sum_{r=1}^n(b_r-a_r)^2}\right).
\]
\end{theorem}

\begin{exercisechunk}
\item \textbf{Relative addressing with RoPE.}
\label[exercise]{ex:rope}
Can a query request a value at a specified relative offset while leaving the
source keys unchanged? We will construct such queries, identify collisions between
addresses, and examine how additional coordinate pairs help distinguish them.

Let $T\ge2$ and $d_v\ge1$ be integers. Position $j\in[T]$ stores an
arbitrary fixed value $v_j\in\mathbb R^{d_v}$. Every source uses the same
unrotated unit key $k\in\mathbb R^2$. For each integer offset $1\le d<T$,
we seek a unit-length query vector $q(d)\in\mathbb R^2$ that retrieves $v_{i-d}$ at every
query position $i\in[T]$ with $i>d$. Reusing this query at different positions
should implement the same addressing rule, independently of the stored values.
As described earlier, RoPE modifies the attention calculation
by applying position-dependent linear
transformations---specifically, rotations---to the query and key vectors before
taking their inner product.

\begin{enumerate}[(a)]
\item \textbf{Encoding the request.}
For a real angle $\phi$, define
\[
R_\phi=\begin{pmatrix}\cos\phi&-\sin\phi\\\sin\phi&\cos\phi\end{pmatrix}.
\]
Fix an \emph{angular frequency} $\omega>0$ (called simply a \emph{frequency}
below): moving forward one position increases the rotation angle by
$\omega$ radians. RoPE compares the
query at position $i$ with the key at position $j$ using the score
\[
s_{ij}=(R_{i\omega}q(d))^\top(R_{j\omega}k),\qquad i,j\in[T].
\]
\begin{enumerate}[label=(\roman*),leftmargin=*,beginpenalty=10000]
\item Show that $R_\alpha^\top R_\beta=R_{\beta-\alpha}$ for real angles $\alpha,\beta$.
\item Find the unique unit vector $q(d)$ so that $s_{i,i-d}=1$ whenever $i>d$, and derive a formula
for $s_{ij}$. Show that no source has a larger score.
\item Verify that shifting
$i$ and $j$ by the same integer amount leaves the score unchanged whenever
both positions remain in $[T]$.
\end{enumerate}
\item \textbf{From scores to retrieval.}
For the scores constructed in (a), let $\varnothing\ne S_i\subseteq[T]$
be the source positions available to query $i$. At positive real scale
$\lambda$, define the attention weights and output by
\[
a_{ij}(\lambda)=\frac{e^{\lambda s_{ij}}}
{\sum_{r\in S_i}e^{\lambda s_{ir}}},\quad j\in S_i,
\qquad
o_i(\lambda)=\sum_{j\in S_i}a_{ij}(\lambda)v_j\in\mathbb R^{d_v}.
\]
Fix $\varepsilon\in(0,1/2)$.
\begin{enumerate}[label=(\roman*),leftmargin=*,beginpenalty=10000]
\item For any $j\in S_i$, show that
\[
\|o_i(\lambda)-v_j\|_2
\le (1-a_{ij}(\lambda))\max_{r\in S_i}\|v_r-v_j\|_2.
\]
Thus $a_{ij}(\lambda)>1-\varepsilon$ gives an $O(\varepsilon)$ retrieval
error for fixed values. More explicitly, for a positive real bound $B$,
show that $\|v_r\|_2\le B$ for all $r\in S_i$ gives error less than
$2B\varepsilon$.
\item Show that guaranteeing error at most $2B\varepsilon$
for every such collection of values requires $a_{ij}(\lambda)\ge1-\varepsilon$.

\item Use $S_i=[i]$ and suppose $0<\omega<2\pi/T$.
Show that $i-d$ is the unique maximizing source at every query position
$i>d$. For each such $i$, show that $\lim_{\lambda\to\infty}o_i(\lambda)=v_{i-d}$,
accounting for the weight at every source position.
\item At any such position, prove that $a_{i,i-d}(\lambda)\ge1-\varepsilon$
requires
\[
\lambda\ge\frac{T^2}{2\pi^2}\log\frac{1-\varepsilon}{\varepsilon}.
\]
\end{enumerate}
Hence, for fixed $\varepsilon$, uniform retrieval over values of norm at most
$B$ requires $\lambda=\Omega(T^2)$ in this construction.
Reducing the frequency makes neighboring scores increasingly close; increasing
$\lambda$ compensates for this shrinking gap. The next part examines collisions
when the frequency is held fixed as the sequence grows, and how a window removes
them.
\item \textbf{Collisions and a window repair.}
For this part, hold $d$ fixed and allow the sequence length $T$ to grow. Choose an integer period $P>d$, set $\omega=2\pi/P$, and use the
query construction from (a) with $S_i=[i]$. A \emph{collision} occurs when
another source has the same score as the requested source $i-d$.
\begin{enumerate}[label=(\roman*),leftmargin=*,beginpenalty=10000]
\item Characterize all colliding source positions.
\item Find the first query position
$i>d$ at which a collision occurs.
\item Evaluate $\lim_{\lambda\to\infty}o_i(\lambda)$ at the first collision
position.
\item Now use the window $S_i=\{j\in[i]:i-j<w\}$, where $w$ is a
positive integer. Find the smallest width that retains $i-d$. For this width,
prove that $i-d$ is the unique maximizing source and that
$\lim_{\lambda\to\infty}o_i(\lambda)=v_{i-d}$ for every $i>d$, however long
the sequence becomes.
\end{enumerate}
\item \textbf{Reliable addressing at a fixed score scale.}
Close scores can require high numerical precision to distinguish, even when
$\lambda$ is large. Can additional coordinate pairs create better-separated
scores and permit a smaller scale? Increasing $\lambda$ is equivalent to
increasing query and key norms. To compare this with increasing width, we fix
$\lambda=1$ and unit query and key norms within each pair.

Fix an integer context length $T\ge2$. Throughout this part, use $S_i=[i]$.
Determine, up to constant factors, how many coordinates are necessary and
sufficient to give the requested
source attention weight at least $3/4$, simultaneously for every integer
offset $1\le d<T$ and every query position $i\in[T]$ with $i>d$.

Let $m$ be a positive integer, which determines the number of coordinate pairs used.
For each coordinate pair $r\in[m]$, use a
fixed unit key $k_r\in\mathbb R^2$, a frequency $\omega_r\in[0,2\pi)$,
and the unit query $q_r(d)\in\mathbb R^2$ constructed in (a).
Define the combined score as
\[
s_{ij}=\sum_{r=1}^m
(R_{i\omega_r}q_r(d))^\top(R_{j\omega_r}k_r),\qquad i,j\in[T].
\]
\begin{enumerate}[label=(\roman*),leftmargin=*,beginpenalty=10000]
\item \textbf{A necessary dimension.}
Show that all scores lie in $[-m,m]$. Use \cref{ex:bounded-attention} to prove
that, at query position $i=T$, attention weight at least $3/4$ on the requested
source requires
\[
m\ge\tfrac12\log\bigl(3(T-1)\bigr).
\]
\item \textbf{Finding frequencies that separate addresses.}
Prove that, for every integer $m\ge8\log(2T)$, there exist frequencies
$\omega_1,\ldots,\omega_m\in[0,2\pi)$ such that
\begin{equation}
\label{eq:rope-separated-addresses}
\sum_{r=1}^m\cos(\omega_r h)\le m/2
\quad\text{for every integer }h\text{ with }1\le|h|<T.
\end{equation}
\item \textbf{Retrieval at every position and offset.}
Fix frequencies satisfying \cref{eq:rope-separated-addresses}. Show that the requested source has score
$m$ and that every other available source has score at most $m/2$,
simultaneously for all offsets and query positions in the question.
Prove that the total attention weight outside the requested source is at most
$(T-1)e^{-m/2}$. Deduce that
$m=\lceil8\log(2T)\rceil$ suffices for target weight at least $3/4$ with
$\lambda=1$. State the resulting error bound when every value vector has
norm at most a positive real bound $B$. Together with (i), conclude that
$\Theta(\log T)$ coordinate pairs are necessary and sufficient.
\end{enumerate}

\end{enumerate}

\paragraph{Conclusion.}
With suitably chosen frequencies, RoPE supports relative addressing with
logarithmic width and bounded coordinates. Extra coordinate pairs both
distinguish addresses and accumulate
the score advantage needed to concentrate attention
(\cref{ex:bounded-attention}). Keeping each coordinate bounded still allows
the total score to grow: width supplies an alternative to increasing $\lambda$.

Standard RoPE instead uses geometrically spaced frequencies to provide both
rapidly and slowly varying positional features. The existence guarantee in
(d) does not automatically apply to that schedule. Indeed, its frequencies
satisfy $0<\omega_r\le1$, so at displacement $h=1$,
\[
\sum_{r=1}^m\cos\omega_r\ge m\cos1>m/2.
\]
Thus it fails the particular sufficient separation condition in (d)(ii);
this does not rule out retrieval with that schedule.

\begin{hints}
For (d)(i), apply the bound from \cref{ex:bounded-attention} with $T$ sources
and score range $2m$.
For (d)(ii), choose the frequencies independently and uniformly on $[0,2\pi)$.
Compute $\mathbb E[\cos(\omega_r h)]$ for a nonzero integer $h$, then apply
\cref{thm:hoeffding} and a union bound over $h\in[T-1]$.
Use the evenness of cosine to handle negative $h$.
\end{hints}
\end{exercisechunk}
